\subsection{Restricción de escalares}
\label{subsection:0.6.5-fr}

\begin{env}[6.5.1]
\label{0.6.5.1-fr}
Sean $A$ un anillo, $\varphi:A\to B$ un homomorfismo de anillos que hace
de $B$ una $A$-álgebra. Supongamos que existe un $B$-módulo $N$ que es un
$A$-módulo \emph{fielmente plano}. Entonces, para todo $A$-módulo $M$,
el homomorfismo $x\to 1\otimes x$ de $M$ en
$B\otimes_A M=M_{(B)}$ es \emph{inyectivo}. En particular, $\varphi$ es
inyectivo; para todo ideal $\mathfrak{a}$ de $A$, se tiene
$\varphi^{-1}(\mathfrak{a}B)=\mathfrak{a}$; para todo ideal maximal (resp.
primo) $\mathfrak{m}$ de $A$, existe un ideal maximal (resp. primo)
$\mathfrak{n}$ de $B$ tal que $\varphi^{-1}(\mathfrak{n})=\mathfrak{m}$.
\end{env}

\begin{env}[6.5.2]
\label{0.6.5.2-fr}
Cuando se cumple la condición de (6.5.1), se identifica $A$ con un
subanillo de $B$ por $\varphi$ y, más generalmente, para todo $A$-módulo
$M$, se identifica $M$ con un sub-$A$-módulo de $M_{(B)}$. Se observará que si
$B$ es entonces \emph{noetheriano}, lo mismo ocurre con $A$, pues la aplicación
$\mathfrak{a}\to\mathfrak{a}B$ es una inyección creciente del
conjunto de los ideales de $A$ en el conjunto de los ideales de $B$;
la existencia de una sucesión infinita estrictamente creciente de ideales de $A$
implicaría, por tanto, la existencia de una sucesión análoga de ideales de $B$.
\end{env}

